% ============================================================
% Appendix D: YAML Test Cases and VHDL Testbenches
% ============================================================
\chapter{YAML Test Cases and VHDL Testbenches}
\label{app:test_cases}

This appendix provides illustrative examples of the key test definition files and simulation scripts used during the verification and validation process, as detailed in \Cref{chap:testing}.

% ------------------------------------------------------------
\section{YAML Test Case Definitions for Fixture-Based Testing}
\label{sec:app_yaml_test_cases}

System-level scenario testing (\Cref{sec:system_scenario_testing}) was driven by sequences defined in \gls{yaml} format and executed via the test fixture~\cite{dolata2023}. The structure allows for the definition of input pin states, expected outputs, and tester interactions.

\subsection*{Example Test Case: Geber 3 Fault Handling}
The following excerpt from \texttt{tests.yml} illustrates the two-step verification sequence for a \enquote{Geber 3} fault:

\begin{verbatim}
# Test group for Geber 3 fault simulation
- name: Geber 3 - I
  description: >-
    Test fur den Geber. Bitte 3 Sekunden warten,
    bis der Fehler angezeigt wird. Danach STW
    gedruckt halten und Fehlersequenz uberprufen.
  pin_sets: {18z: 1, 24d: 0} # Simulate fault
  expected_err_code: "31"    # Expect immediate err 31
  expected_err_seq: "03 31 09" # Expect seq on STW
  del: False                 # No 'Loesch' prompt

- name: Geber 3 - II
  description: STW gedruckt halten und Fehlersequenz uberprufen.
  pin_sets: {}               # Maintain fault state
  expected_err_code: "95"    # Expect '95' on STW
  expected_err_seq: "03 31 09" # Re-verify seq
  del: True                  # Prompt 'Loesch'
\end{verbatim}

% ------------------------------------------------------------
\section{VHDL Snippet for Simulation-Based Verification}
\label{sec:app_vhdl_testbenches}

Component and interface verification (\Cref{sec:fpga_interface_testing}) utilized deterministic \gls{vhdl} testbenches. The following code snippet from the top-level FPGA module (`atop.vhd`) illustrates the control logic for the \gls{nvram}, including the critical bug-fix for the \texttt{NVRAM\_ble} (active-low enable) signal, which was a key finding of the practical debugging phase~\cite{savchuk2025}.

\subsection*{Example Snippet: NVRAM Control Logic}
This code block is responsible for asserting the correct control signals (chip enable, output enable, write enable) for the non-volatile memory.

\begin{verbatim}
-- VHDL snippet from NVRAM control process in atop.vhd

-- Active-low enable pin for NVRAM. Was incorrectly set to '1' in
-- an earlier version, which disabled the chip.
NVRAM_ble <= '0'; -- FIX: Set permanently low to enable the component.

NVRAM_oe_s <= rd_n_os;
NVRAM_ce_s <= not ((not (rd_n_os or p27)) or (not (p27 or wr_n_os)));
NVRAM_ce <= NVRAM_ce_s;
NVRAM_we <= wr_n_os;

-- Data bus direction control
dmem_data_is <= NVRAM_dq;

-- Other control signals for the memory controller state machine
NVRAM_ctrl(2) <= p2_os(7);
NVRAM_ctrl(1) <= wr_n_os;
NVRAM_ctrl(0) <= NVRAM_ce_s;
\end{verbatim}